% !TEX program = pdflatex
% ════════════════════════════════════════════════════════════════════════════
%  UrbanFlow CDMX / VialAI — Artículo del proyecto integrador
%  Diplomado en Ciencia de Datos · FES Acatlán, UNAM · Gen 30 · 2026
%  Autor: Carlos Armando López Encino
%  Profesor: Mtro. Fernando Barranco Rodríguez
%  Versión v3 — abril 2026
% ════════════════════════════════════════════════════════════════════════════

\documentclass[11pt,a4paper]{article}

% ───── Codificación y lenguaje ─────────────────────────────────
\usepackage[utf8]{inputenc}
\usepackage[T1]{fontenc}
\usepackage[spanish,es-tabla,es-noindentfirst]{babel}
\usepackage[round,authoryear]{natbib}
\usepackage{lmodern}
\usepackage{microtype}

% ───── Matemáticas ─────────────────────────────────────────────
\usepackage{amsmath,amssymb,amsthm,mathtools}
\usepackage{bm}
\usepackage{siunitx}
\sisetup{output-decimal-marker={.},group-separator={\,}}

% ───── Geometría y espaciado ───────────────────────────────────
\usepackage[a4paper,margin=2.5cm,headheight=15pt]{geometry}
\usepackage{setspace}
\setstretch{1.15}
\usepackage{parskip}

% ───── Color e identidad VialAI ────────────────────────────────
\usepackage{xcolor}
\definecolor{vialNavyDeep}{HTML}{0B1F3A}
\definecolor{vialNavy}{HTML}{0F2E4D}
\definecolor{vialTeal}{HTML}{2EBFA6}
\definecolor{vialAqua}{HTML}{3BA3BD}
\definecolor{vialGreen}{HTML}{A3D88F}
\definecolor{vialAmber}{HTML}{F5B836}
\definecolor{vialText}{HTML}{1F2937}
\definecolor{vialMuted}{HTML}{4A5B6D}
\definecolor{vialBgSoft}{HTML}{F3F6FA}

% ───── Gráficos y figuras ──────────────────────────────────────
\usepackage{graphicx}
% Rutas donde buscar imágenes (el . es la carpeta del .tex)
\graphicspath{{./}{graficas_proyecto/imagenes_12/}}
\usepackage{float}
\usepackage{subcaption}
\usepackage{wrapfig}
\usepackage{tikz}
\usetikzlibrary{arrows.meta,positioning,shapes.geometric,shadows.blur,
                calc,fit,backgrounds,decorations.pathreplacing,
                shapes.symbols,shapes.multipart,chains,scopes}

% ───── Tablas ──────────────────────────────────────────────────
\usepackage{booktabs}
\usepackage{array}
\usepackage{tabularx}
\usepackage{multirow}
\usepackage{longtable}
\usepackage{colortbl}

% ───── Tipografía de secciones ─────────────────────────────────
\usepackage{titlesec}
\titleformat{\section}
  {\color{vialNavy}\normalfont\Large\bfseries}{\thesection.}{0.5em}{}
\titleformat{\subsection}
  {\color{vialNavy}\normalfont\large\bfseries}{\thesubsection}{0.5em}{}
\titleformat{\subsubsection}
  {\color{vialTeal}\normalfont\normalsize\bfseries}{\thesubsubsection}{0.5em}{}

% ───── Encabezados y pies ──────────────────────────────────────
\usepackage{fancyhdr}
\pagestyle{fancy}
\fancyhf{}
\fancyhead[L]{\color{vialMuted}\small VialAI --- Predicción estocástica ZMVM}
\fancyhead[R]{\color{vialMuted}\small\thepage}
\renewcommand{\headrulewidth}{0.4pt}
\renewcommand{\headrule}{\color{vialTeal}\hrule width\headwidth}

% ───── Cajas de callout ────────────────────────────────────────
\usepackage[most]{tcolorbox}

\newtcolorbox{tipbox}[1][]{%
  enhanced, breakable,
  colback=vialBgSoft, colframe=vialTeal, coltitle=vialNavy,
  fonttitle=\bfseries, title=#1,
  boxrule=0.8pt, arc=3mm, left=4mm, right=4mm, top=2mm, bottom=2mm,
  drop shadow=vialMuted,
}

\newtcolorbox{keybox}[1][]{%
  enhanced, breakable,
  colback=vialNavy!5, colframe=vialNavy, coltitle=white,
  fonttitle=\bfseries, title=#1,
  boxrule=1pt, arc=3mm, left=4mm, right=4mm, top=2mm, bottom=2mm,
}

\newtcolorbox{casebox}[1][]{%
  enhanced, breakable,
  colback=vialAmber!10, colframe=vialAmber!80!black, coltitle=vialNavy,
  fonttitle=\bfseries, title=#1,
  boxrule=0.8pt, arc=3mm, left=4mm, right=4mm, top=2mm, bottom=2mm,
}

\newtcolorbox{warnbox}[1][]{%
  enhanced, breakable,
  colback=red!5, colframe=red!60!black, coltitle=white,
  fonttitle=\bfseries, title=#1,
  boxrule=0.8pt, arc=3mm, left=4mm, right=4mm, top=2mm, bottom=2mm,
}

% ───── Hipervínculos ───────────────────────────────────────────
\usepackage[colorlinks=true,
            linkcolor=vialNavy,
            citecolor=vialTeal,
            urlcolor=vialAqua]{hyperref}

% ───── Código ──────────────────────────────────────────────────
\usepackage{listings}
\lstdefinestyle{vialcode}{
  backgroundcolor=\color{vialBgSoft},
  basicstyle=\ttfamily\footnotesize\color{vialText},
  keywordstyle=\color{vialNavy}\bfseries,
  stringstyle=\color{vialTeal},
  commentstyle=\color{vialMuted}\itshape,
  frame=single, framesep=5pt, rulecolor=\color{vialTeal},
  numbers=left, numberstyle=\tiny\color{vialMuted},
  breaklines=true, showstringspaces=false, captionpos=b,
}
\lstset{
  literate={á}{{\'a}}1 {é}{{\'e}}1 {í}{{\'i}}1 {ó}{{\'o}}1 {ú}{{\'u}}1
           {Á}{{\'A}}1 {É}{{\'E}}1 {Í}{{\'I}}1 {Ó}{{\'O}}1 {Ú}{{\'U}}1
           {ñ}{{\~n}}1 {Ñ}{{\~N}}1 {¿}{?`}1 {¡}{!`}1
}

% ───── Comandos propios ────────────────────────────────────────
\newcommand{\VialAI}{\textbf{\textcolor{vialTeal}{Vial}\textcolor{vialAmber}{AI}}}
\newcommand{\ZMVM}{\textsc{zmvm}}
\newcommand{\CDMX}{\textsc{cdmx}}
\newcommand{\AICM}{\textsc{aicm}}
\newcommand{\MAPE}{\text{MAPE}}
\newcommand{\Prob}{\mathbb{P}}
\newcommand{\Exp}{\mathbb{E}}
\newcommand{\Var}{\mathrm{Var}}

\begin{document}

% ═══════════════════════ PORTADA ═══════════════════════════════
\begin{titlepage}
\centering
\vspace*{1cm}

{\color{vialMuted}\large Diplomado en Ciencia de Datos \textbullet\ Generación 30}\\[0.3em]
{\color{vialMuted}\small FES Acatlán \textbullet\ UNAM \textbullet\ 2026}\\[1.2cm]


\includegraphics[width=0.28\textwidth]{Logo_VialAI.png}\\[1em]

{\color{vialTeal}\Huge\bfseries UrbanFlow \CDMX}\\[0.5em]
{\color{vialMuted}\Large\itshape Inteligencia para la movilidad urbana}\\[1.3cm]

{\color{vialNavy}\Large\bfseries Predicción estocástica de tiempos de viaje\\
en la Zona Metropolitana del Valle de México:}\\[0.4em]
{\color{vialNavy}\large un sistema multimodal basado en cadenas de Markov,\\
simulación Monte Carlo y agentes conversacionales}\\[2cm]

{\large\textbf{Carlos Armando López Encino}}\\
{\small\color{vialMuted}cale.delta@gmail.com}\\[0.8em]
{\small Profesor: Mtro. Fernando Barranco Rodríguez}\\[2.3cm]

{\color{vialMuted}\small Proyecto integrador \textbullet\ 18 de abril de 2026}\\[0.4em]
{\color{vialAqua}\small\url{https://github.com/caledelta/UrbanFlow_CDMX}}

\vfill
{\color{vialMuted}\footnotesize
Este documento forma parte del proyecto integrador final del Diplomado
en Ciencia de Datos. Todo el código, datos de calibración y resultados
son reproducibles desde el repositorio público.}
\end{titlepage}

% ---------- PÁGINA EN BLANCO ----------
\newpage
\mbox{} % Genera un elemento invisible para que la página no se ignore
\newpage

% ═══════════════════════ RESUMEN ═══════════════════════════════
\begin{abstract}
\noindent
La Zona Metropolitana del Valle de México concentra a
\num{21800000} habitantes que pierden en promedio \num{152} horas al año
atascados en el tráfico, según el índice global de congestión de TomTom.
Los sistemas actuales de estimación de tiempos de viaje devuelven un
único número determinista que oculta la variabilidad real del tráfico
urbano. Este trabajo presenta \VialAI\ (UrbanFlow \CDMX), un sistema
de predicción \emph{probabilística} que combina cadenas de Markov
calibradas con datos reales del C5 \CDMX\ y TomTom Traffic Stats con
simulación Monte Carlo vectorizada ($N=\num{10000}$ trayectorias),
expuesto al usuario final mediante una interfaz multimodal que acepta
entrada por voz en español mexicano y razona con un agente basado en
Claude 4.6 dotado de seis herramientas especializadas. La arquitectura
está diseñada con \emph{degradación graceful} en tres capas (transcripción
de voz, compresión semántica con LLM y síntesis de voz) para garantizar
resiliencia ante fallas de infraestructura externa. Los resultados
preliminares sobre una muestra de \num{30} rutas de campo en la \ZMVM\
reportan un MAPE del \SI{11.9}{\percent} en la mediana predicha
(contra \SI{19.0}{\percent} del ETA de TomTom, una ventaja de
\SI{7.1}{pp}) y una cobertura empírica del \SI{93.3}{\percent} para la
banda nominal P10--P90. Se discute el nicho B2B de conductores de
última milla y se proyecta una campaña de validación cuantitativa con
$N \geq \num{500}$ para la siguiente iteración.

\vspace{0.4em}
\noindent\textbf{Palabras clave:} predicción estocástica, cadenas de
Markov, Monte Carlo, tráfico urbano, \ZMVM, agentes conversacionales,
interfaz por voz, Streamlit, Claude, degradación \emph{graceful}.
\end{abstract}

\begin{abstract}\selectlanguage{english}
\noindent
The Mexico City Metropolitan Area is home to \num{21.8} million
inhabitants who lose an average of \num{152} hours per year to
traffic congestion, according to TomTom's global congestion index.
Current travel-time estimation systems return a single deterministic
number that hides the intrinsic variability of urban traffic. This
work presents \textbf{VialAI} (UrbanFlow CDMX), a \emph{probabilistic}
prediction system combining Markov chains calibrated with real data
from Mexico City's C5 emergency service and TomTom Traffic Stats
with vectorised Monte Carlo simulation ($N=\num{10000}$ trajectories),
exposed to end users via a multimodal interface that accepts voice
input in Mexican Spanish and reasons through a Claude 4.6-based agent
equipped with six specialised tools. The architecture features
\emph{graceful degradation} across three layers (speech-to-text,
semantic-compression LLM, and text-to-speech) to ensure resilience
against external infrastructure failures. Preliminary results on a
\num{30}-route field sample in the metropolitan area report an
\SI{11.9}{\percent} MAPE on the predicted median (versus
\SI{19.0}{\percent} for TomTom's ETA, a \SI{7.1}{pp} advantage) and an
empirical coverage of \SI{93.3}{\percent} for the nominal P10--P90
band. We discuss the B2B niche of last-mile drivers and outline a
quantitative validation campaign with $N \geq \num{500}$ for the
next iteration.

\vspace{0.4em}
\noindent\textbf{Keywords:} stochastic prediction, Markov chains,
Monte Carlo, urban traffic, Mexico City, conversational agents,
voice interface, Streamlit, Claude, graceful degradation.
\selectlanguage{spanish}\end{abstract}

\vspace{1em}
\noindent\hrulefill

% ---------- PÁGINA EN BLANCO ----------
\newpage
\mbox{} % Genera un elemento invisible para que la página no se ignore

\newpage
\tableofcontents
\newpage

% ═══════════════════════ §1 INTRODUCCIÓN ═══════════════════════
\section{Introducción}
\label{sec:intro}

\subsection{El problema: \CDMX\ como laboratorio de congestión urbana}

La Ciudad de México y su zona metropolitana constituyen uno de los
laboratorios más demandantes del mundo para el estudio de la
movilidad urbana. Con una densidad de \num{5966} habitantes por
kilómetro cuadrado en la \CDMX\ propiamente dicha
\citep{inegi2024densidad}, una flota vehicular que supera los
\num{5800000} unidades registradas \citep{semovi2024parque} y una
infraestructura vial diseñada originalmente para la mitad de esa
demanda, el resultado es predecible y persistente: la \ZMVM\ ocupa
recurrentemente los primeros diez lugares del ranking global de
ciudades más congestionadas \citep{tomtom2024index}.

La Figura~\ref{fig:ranking} resume el posicionamiento relativo de la
\CDMX\ contra otras megalópolis comparables.

\begin{figure}[H]
\centering
\includegraphics[width=0.88\textwidth]{fig1_congestion_ranking.png}
\caption{Ranking de congestión urbana: tiempo promedio perdido al
año por conductor en horas (fuente: TomTom Traffic Index, edición
\num{2024}). La \CDMX\ pierde \num{152} horas anuales por
conductor, equivalente a \num{6.3} días completos.}
\label{fig:ranking}
\end{figure}

Detrás de ese número agregado se esconde una realidad más problemática
para quien tiene que tomar decisiones concretas de ruta y horario.
La congestión no es simplemente alta: es \textbf{altamente variable}.
Un trayecto Polanco--\AICM\ de \num{13} kilómetros puede tomar
\num{18} minutos a las \num{6} de la mañana de un domingo y
\num{72} minutos a las \num{19} de un jueves lluvioso. Esa
volatilidad hace que las herramientas tradicionales de estimación
de tiempo de viaje fallen en la pregunta que realmente importa:
no ``\emph{¿cuánto tardaré?}'' sino ``\emph{¿cuánto puedo tardar
en el peor caso razonable?}''

\subsection{¿Por qué inteligencia artificial, aprendizaje automático,
cadenas de Markov y Monte Carlo \emph{juntos}?}
\label{sec:por-que-juntos}

La predicción de tiempos de viaje se ha abordado históricamente desde
tres familias de técnicas: modelos estadísticos clásicos (regresión,
series temporales), aprendizaje supervisado sobre trayectorias
históricas (GBM, redes neuronales recurrentes) y métodos de
simulación discreta para sistemas de tráfico. Cada familia captura
una pieza del fenómeno, pero ninguna aislada responde bien a la
pregunta del intervalo de incertidumbre contextualizado.

\VialAI\ combina deliberadamente cuatro paradigmas complementarios:

\begin{keybox}[Las cuatro piezas del razonamiento de VialAI]
\textbf{\textcolor{vialTeal}{Cadenas de Markov}} capturan la
estructura dinámica del tráfico urbano modelándolo como un proceso
estocástico de tres estados (fluido, moderado, congestionado) con
matrices de transición calibradas por hora del día y día de la
semana. Proporcionan la física probabilística del sistema.\\[0.3em]
\textbf{\textcolor{vialTeal}{Simulación Monte Carlo}} permite
propagar la incertidumbre de cada transición a lo largo del trayecto
completo. Al generar $N = \num{10000}$ trayectorias independientes
y calcular percentiles, obtenemos una distribución empírica del
tiempo de viaje en vez de un valor puntual.\\[0.3em]
\textbf{\textcolor{vialAmber}{Aprendizaje automático}} calibra las
probabilidades de transición desde datos reales: velocidades
históricas de TomTom Traffic Stats, incidentes del C5 \CDMX\ y
condiciones meteorológicas de OpenWeatherMap. Esto convierte el
modelo teórico en un modelo ajustado a la realidad operativa de la
\ZMVM.\\[0.3em]
\textbf{\textcolor{vialAmber}{Inteligencia artificial generativa}}
(Claude Opus 4.6 y Haiku 4.5) actúa como capa de interpretación
semántica: traduce lenguaje natural del conductor a invocaciones
estructuradas de las seis herramientas del sistema, inyecta contexto
de lugares personalizados y comprime la respuesta matemática en una
oración hablable de menos de \num{30} palabras.
\end{keybox}

El valor emerge de la \emph{composición}. Las cadenas de Markov sin
calibración ML producirían predicciones académicamente correctas
pero operativamente inútiles. El ML sin estructura estocástica
produciría puntos estimados sin banda de incertidumbre. Monte Carlo
sin agente produciría intervalos rigurosos que el conductor no puede
consumir mientras maneja. Y un agente conversacional sin motor
matemático sería un chatbot que alucina números. Juntas, estas
cuatro capas conforman un sistema que es \emph{matemáticamente
riguroso}, \emph{empíricamente calibrado}, \emph{operativamente
accionable} y \emph{accesible al usuario final}.

\subsubsection*{Un caso ilustrativo: Miguel, repartidor de última
milla}

\begin{casebox}[El conductor de última milla real]
\small\noindent
Son las 18:40 de un miércoles lluvioso. Miguel acaba de
entregar un pedido en Polanco y acepta uno nuevo con destino en
Iztapalapa. La aplicación le muestra un ETA de \num{42} minutos.
Miguel sabe por experiencia que ese número es engañoso: a veces ha
tardado \num{35}, a veces \num{70}. Cuando acepta, apuesta su bono
por entregas a tiempo contra una distribución que no conoce.
\VialAI\ le responde con voz \emph{``Son aproximadamente
\num{12.4} kilómetros, tardarás unos \num{58} minutos, pero la banda
realista va de \num{44} a \num{78} con lluvia moderada en
Circuito Interior''}. Esa frase, que se genera en menos de
\num{4} segundos desde que Miguel termina de hablar, contiene el
\emph{valor} agregado del sistema: la mediana contextualizada
(P50), la banda P10--P90 que enmarca la incertidumbre y la razón
dominante (\emph{lluvia en Circuito}) que lo explica.
\end{casebox}

Miguel es el usuario canónico de este trabajo. No es un ingeniero
de software, no va a tocar un botón mientras maneja, y cada minuto
de error en la estimación le cuesta dinero real o seguridad
personal. Diseñar para Miguel ha sido la restricción de diseño
dominante.

\subsection{Hipótesis, objetivos y contribuciones}

\textbf{Hipótesis.} Es posible construir un sistema de predicción
de tiempos de viaje para la \ZMVM\ que (i)~reporte un intervalo
de confianza calibrado empíricamente, (ii)~integre fuentes
heterogéneas de datos en tiempo real y cuasi-tiempo-real,
(iii)~sea consultable por voz natural en español mexicano, y
(iv)~opere en presencia de fallas de servicios externos, todo con
latencia inferior a \num{5} segundos por consulta y desplegable
en hardware de gama media.

\textbf{Objetivos.}

\begin{enumerate}
\item Formalizar el problema como un proceso estocástico de tres
estados calibrado con datos operativos.
\item Implementar simulación Monte Carlo vectorizada con
$N \geq \num{10000}$ y evaluar convergencia.
\item Construir un pipeline de ingesta para cinco fuentes de datos
heterogéneas.
\item Diseñar un agente conversacional con seis herramientas que
traduzca lenguaje natural a llamadas estructuradas.
\item Exponer la funcionalidad mediante una interfaz multimodal
(voz, texto, clic) con arquitectura de tres capas con
degradación \emph{graceful}.
\item Validar empíricamente el sistema contra el ETA comercial de
referencia (TomTom Routing).
\end{enumerate}

\textbf{Contribuciones.} Este trabajo aporta: (1)~una
implementación abierta y reproducible del pipeline completo;
(2)~una calibración empírica de la matriz de transición para tres
corredores de la \ZMVM; (3)~evidencia preliminar ($N=30$) de una
mejora de \SI{7.1}{pp} en MAPE respecto al estado del arte
comercial; (4)~una arquitectura de voz de tres capas con
degradación \emph{graceful} documentada como patrón reutilizable;
y (5)~un marco de referencia para monetización B2B orientado a
servicios de última milla.

\subsection{Organización del documento}

La Sección~\ref{sec:identidad} describe la identidad y estructura
del proyecto, incluyendo el pipeline visual completo y el stack
tecnológico. La Sección~\ref{sec:arte} revisa el estado del arte
y posiciona las contribuciones. La Sección~\ref{sec:nicho}
caracteriza el nicho de mercado y el perfil del usuario objetivo.
La Sección~\ref{sec:datos} detalla las fuentes de datos y sus
flujos de ingesta. Las Secciones~\ref{sec:modelo} y
\ref{sec:perturbaciones} desarrollan el motor estocástico y el
tratamiento de perturbaciones. La Sección~\ref{sec:agente}
presenta el agente conversacional. La Sección~\ref{sec:voz}
expone la arquitectura de voz multimodal. La
Sección~\ref{sec:app} describe la aplicación web como producto
entregable. La Sección~\ref{sec:validacion} reporta los
resultados empíricos. La Sección~\ref{sec:futuro} discute
limitaciones y trabajo futuro. La Sección~\ref{sec:conclusiones}
cierra con conclusiones.

% ═══════════════════════ §2 IDENTIDAD ══════════════════════════
\section{Identidad y estructura del proyecto}
\label{sec:identidad}

\subsection{Marca e identidad visual}

\VialAI\ nace como un proyecto de Diplomado pero está diseñado con
la ambición de sobrevivir más allá del entregable académico. Su
identidad visual, que aparece en el logotipo de la portada, está
construida alrededor de tres ideas: (i)~la \emph{intersección}
como metáfora de la ruta que converge desde múltiples orígenes
posibles (el símbolo en forma de $\vee$); (ii)~la \emph{movilidad
inteligente} como promesa de producto (el gradiente del verde-agua
al verde-lima); y (iii)~la \emph{confianza técnica} como soporte
institucional (el fondo azul marino profundo).

\begin{figure}[H]
\centering
\begin{tikzpicture}
% Fondo
\fill[vialNavyDeep, rounded corners=4pt] (0,0) rectangle (15,2.6);
% Muestras de color con nombre
\foreach \clr/\name/\hex/\x in {%
  vialNavyDeep/Fondo profundo/0B1F3A/0.3,
  vialNavy/Navy principal/0F2E4D/2.8,
  vialTeal/Verde-agua/2EBFA6/5.3,
  vialAqua/Azul agua/3BA3BD/7.8,
  vialGreen/Verde suave/A3D88F/10.3,
  vialAmber/Ambar marca/F5B836/12.8%
}{%
  \fill[\clr, rounded corners=3pt] (\x,0.4) rectangle (\x+2.0,2.0);
  \node[white, font=\scriptsize\bfseries, align=center] at (\x+1.0,1.4) {\name};
  \node[white!80, font=\tiny\ttfamily] at (\x+1.0,0.8) {\#\hex};
}
\end{tikzpicture}
\caption{Paleta cromática oficial de \VialAI. Los tres tonos fríos
(navy, teal, aqua) se emplean en elementos estructurales y de
datos. Los dos tonos cálidos (verde suave y ámbar) se reservan
para destacar acción del usuario y estados positivos del sistema.}
\label{fig:paleta}
\end{figure}

\subsection{Pipeline completo del sistema}

La Figura~\ref{fig:pipeline} sintetiza la arquitectura end-to-end.
El pipeline se lee de izquierda (entrada del usuario) a derecha
(respuesta accionable) y distingue cuatro capas funcionales
claramente delimitadas.

\begin{figure}[H]
\centering
\begin{tikzpicture}[
  node distance=0.7cm and 1.0cm,
  layer/.style={rectangle, rounded corners=3pt, draw=vialNavy,
                fill=vialNavy!6, minimum width=2.7cm, minimum height=1.2cm,
                align=center, font=\small, drop shadow=vialMuted},
  voice/.style={layer, fill=vialTeal!15, draw=vialTeal},
  data/.style={layer, fill=vialAqua!15, draw=vialAqua},
  core/.style={layer, fill=vialAmber!20, draw=vialAmber!80!black},
  salida/.style={layer, fill=vialGreen!25, draw=vialGreen!60!black},
  arr/.style={-Latex, line width=0.8pt, draw=vialNavy},
  dashedarr/.style={-Latex, line width=0.6pt, draw=vialMuted, dashed},
]
% Col 1: Entrada
\node[voice] (voz) {\textbf{Voz}\\ \footnotesize faster-whisper};
\node[voice, below=of voz] (chat) {\textbf{Chat}\\ \footnotesize texto libre};
\node[voice, below=of chat] (manual) {\textbf{Manual}\\ \footnotesize mapa + selectbox};

% Col 2: Agente
\node[core, right=1.5cm of chat] (agente) {\textbf{Agente}\\ \textbf{VialAI}\\
  \footnotesize Claude 4.6\\ \footnotesize + 6 tools};

% Col 3: Datos
\node[data, right=1.5cm of agente, yshift=1.6cm] (tomtom) {TomTom\\ \footnotesize Traffic + Routing};
\node[data, right=1.5cm of agente] (c5) {C5 \CDMX\\ \footnotesize CKAN API};
\node[data, right=1.5cm of agente, yshift=-1.6cm] (owm) {OpenWeather\\ \footnotesize clima actual};

% Col 4: Motor
\node[core, right=1.5cm of c5] (motor) {\textbf{Motor}\\ \textbf{estocástico}\\
  \footnotesize Markov + MC\\ $N=\num{10000}$};

% Col 5: Salida
\node[salida, right=1.5cm of motor, yshift=1cm] (ui) {\textbf{UI web}\\ \footnotesize Streamlit + Folium};
\node[salida, right=1.5cm of motor, yshift=-1cm] (tts) {\textbf{Voz saliente}\\ \footnotesize TTS + Haiku};

% Flechas entrada -> agente
\draw[arr] (voz.east) -- (agente.west);
\draw[arr] (chat.east) -- (agente.west);
\draw[arr] (manual.east) -- ($(agente.west)+(0,-0.2)$);

% Flechas agente -> datos
\draw[arr] (agente.east) -- (tomtom.west);
\draw[arr] (agente.east) -- (c5.west);
\draw[arr] (agente.east) -- (owm.west);

% Flechas datos -> motor
\draw[arr] (tomtom.east) -- (motor.west);
\draw[arr] (c5.east) -- (motor.west);
\draw[arr] (owm.east) -- (motor.west);

% Flechas motor -> salida
\draw[arr] (motor.east) -- (ui.west);
\draw[arr] (motor.east) -- (tts.west);

% Retorno de TTS al usuario
\draw[dashedarr] (tts.east) to[bend right=60]
  node[right=5pt, font=\scriptsize, text=vialMuted] {}
  ($(voz.east)+(8,0)$) -- (voz.east);
\end{tikzpicture}
\caption{Pipeline funcional de \VialAI. Entrada multimodal
(izquierda) se traduce en llamadas estructuradas por el agente
(centro), que a su vez consulta fuentes heterogéneas y alimenta
el motor estocástico (centro-derecha) cuyo resultado se expone al
usuario en dos canales: visual vía Streamlit y auditivo vía TTS
resumido con Claude Haiku (derecha). La flecha punteada indica el
loop auditivo de vuelta al conductor.}
\label{fig:pipeline}
\end{figure}

\subsection{Stack tecnológico y herramientas}
\label{sec:stack}

La selección de herramientas siguió tres criterios: (i)~madurez
productiva comprobada, (ii)~reproducibilidad desde el repositorio
público, y (iii)~coherencia con el ecosistema del Diplomado.
La Tabla~\ref{tab:stack} resume el stack completo.

\begin{table}[H]
\centering
\small
\renewcommand{\arraystretch}{1.2}
\begin{tabularx}{\linewidth}{@{}l l X@{}}
\toprule
\rowcolor{vialNavy!10}
\textbf{Capa} & \textbf{Herramienta} & \textbf{Rol en el proyecto} \\
\midrule
Lenguaje base & Python \num{3.12} & Todo el motor estocástico, agente,
ingesta y orquestación. Elegido por madurez del ecosistema
\texttt{numpy}/\texttt{scipy} y por ser el estándar del Diplomado.\\
Exploración & Google Colab & Prototipado inicial del modelo Markov
y visualización de distribuciones durante la fase de diseño. El
notebook de referencia está versionado en \texttt{notebooks/}.\\
UI web & Streamlit \num{1.56} & Interfaz principal del usuario final.
Elegido por rapidez de iteración y por ofrecer hot-reload sobre
código Python sin separar front de back.\\
Mapa & Folium + Leaflet & Visualización geoespacial interactiva.
Permite renderizar rutas reales de TomTom Routing como polilíneas
coloreadas y marcadores de lugares guardados con iconografía
contextual.\\
Agente & Claude Opus \num{4.6} & Motor de razonamiento del agente
VialAI con \emph{tool use} en bucle cerrado. Elegido por capacidad
superior en tareas de orquestación de herramientas y generación en
español mexicano.\\
Compresión & Claude Haiku \num{4.5} & Generación del resumen
hablable (máx.\ \num{30} palabras) a partir de la respuesta larga
con markdown del agente. Seleccionado por latencia ($<$\SI{1}{s})
y costo ($\sim$\SI{0.0001}USD por consulta).\\
STT primario & faster-whisper & Transcripción local y gratuita
de audio a texto. Corre en CPU sin dependencia de APIs externas.\\
STT respaldo & Whisper-1 (OpenAI) & Alternativa cloud con mayor
calidad de transcripción para casos donde el usuario active la
ruta premium.\\
TTS primario & OpenAI tts-1 (voz \emph{nova}) & Síntesis de voz
natural para la respuesta hablada. Calidad de demo.\\
TTS respaldo & pyttsx3 (SAPI5) & Síntesis local offline para
resiliencia ante fallas del TTS cloud.\\
Testing & pytest \num{7.4} & \num{720}+ tests automatizados cubren
motor, ingesta, agente, persistencia y voz.\\
Control de versiones & Git + GitHub & Historial público completo
del desarrollo. Cada hito está referenciado por hash de commit.\\
Asistencia IA & Claude Code & Herramienta de pair programming
usada durante el desarrollo para refactors estructurales y
generación de tests. Todas las decisiones de diseño son del autor;
la asistencia se empleó principalmente para acelerar iteración sin
sacrificar revisión crítica.\\
\bottomrule
\end{tabularx}
\caption{Stack tecnológico de \VialAI. La selección prioriza
herramientas abiertas, reproducibles y con alternativas de
respaldo cuando el componente depende de servicios externos.}
\label{tab:stack}
\end{table}

\begin{tipbox}[Principio de selección de herramientas]
Cada componente con dependencia externa (voz, LLM, rutas) tiene
una alternativa local o de respaldo implementada y probada. Esta
decisión de diseño es la base de la arquitectura de
\emph{degradación graceful} que se detalla en la
Sección~\ref{sec:voz}.
\end{tipbox}

% ═══════════════════════ §3 ESTADO DEL ARTE ════════════════════
\section{Estado del arte y diferenciación}
\label{sec:arte}

La predicción de tiempos de viaje urbanos es uno de los problemas
más estudiados en ingeniería del tráfico y ciencia de datos
geoespaciales. Esta sección posiciona a \VialAI\ contra tres
familias de soluciones preexistentes y explicita en qué ejes se
diferencia.

\subsection{Familia 1: ETAs comerciales deterministas}

Las aplicaciones de navegación masivas (Google Maps, Waze, Apple
Maps, TomTom) y las \textsc{api} de routing que las respaldan
entregan al usuario final un \emph{único número} como tiempo
estimado de arribo. Internamente, esos servicios sí modelan
distribuciones de velocidad segmento a segmento, pero el producto
expuesto colapsa la incertidumbre a la mediana o a un promedio
ponderado \citep{google2023eta,tomtom2024routing}.

Para un usuario casual esto es suficiente. Para un conductor
profesional con incentivos económicos atados a entregas a tiempo,
la ocultación de la varianza es un problema estructural: el ETA
de \num{42} minutos que Google muestra puede ocultar una
distribución con $\sigma = 15$ minutos, lo que se traduce en una
probabilidad no trivial de llegar entre \num{27} y \num{57}
minutos. Esa diferencia define si el conductor cobra bono o no.

\subsection{Familia 2: Modelos académicos de predicción probabilística}

En la literatura académica, la predicción de tiempos de viaje con
cuantificación de incertidumbre se ha abordado con regresión de
cuantiles \citep{koenker2005quantile}, redes neuronales bayesianas
\citep{gal2016dropout}, modelos de mezcla gaussiana
\citep{jenelius2013travel} y, más recientemente, transformers
espaciotemporales con salida probabilística
\citep{li2023stformer}.

Estos trabajos son técnicamente sólidos pero adolecen de tres
limitaciones operativas para el caso de uso que \VialAI\ aborda:

\begin{enumerate}
\item Requieren entrenamiento supervisado sobre conjuntos masivos
de trayectorias \textsc{gps} que no están disponibles públicamente
para la \ZMVM.
\item Su interpretación de la incertidumbre es estadística pero
no \emph{explicativa}: un intervalo amplio no viene acompañado de
la causa dominante (\emph{llueve}, \emph{hay un evento deportivo}).
\item Carecen de una capa de interfaz natural al usuario final;
operan como modelos de \emph{backend} que otros ingenieros deben
integrar.
\end{enumerate}

\subsection{Familia 3: Sistemas basados en simulación microscópica}

Los simuladores de tráfico clásicos (\textsc{sumo}, Vissim,
Aimsun) modelan explícitamente cada vehículo y sus interacciones
\citep{lopez2018sumo}. Son extraordinariamente precisos para
estudios de planeación urbana de largo plazo pero requieren
modelado detallado de la red vial completa, calibración intensiva
con aforos vehiculares, y tiempos de ejecución que los hacen
inviables para predicciones en línea con latencia inferior a
\num{5} segundos.

\subsection{Posición de VialAI en el espacio de soluciones}

\VialAI\ ocupa deliberadamente un punto distinto en el espacio
de trade-offs que ninguna de las tres familias anteriores cubre
simultáneamente. La Tabla~\ref{tab:posicionamiento} lo resume.

\begin{table}[H]
\centering
\small
\renewcommand{\arraystretch}{1.3}
\begin{tabularx}{\linewidth}{@{}X c c c c@{}}
\toprule
\rowcolor{vialNavy!10}
\textbf{Eje de diferenciación} & \textbf{ETA comercial} &
\textbf{ML académico} & \textbf{Simuladores} &
\textbf{\VialAI} \\
\midrule
Intervalo probabilístico P10--P90 & \textcolor{vialAmber}{parcial} &
\textcolor{vialGreen}{sí} & \textcolor{vialGreen}{sí} &
\textcolor{vialGreen}{\textbf{sí}} \\
Explicación causal de la varianza & \textcolor{red}{no} &
\textcolor{red}{no} & \textcolor{vialAmber}{parcial} &
\textcolor{vialGreen}{\textbf{sí}} \\
Interfaz natural (voz/chat) & \textcolor{red}{no} &
\textcolor{red}{no} & \textcolor{red}{no} &
\textcolor{vialGreen}{\textbf{sí}} \\
Latencia $<\SI{5}{s}$ por consulta &
\textcolor{vialGreen}{sí} & \textcolor{vialAmber}{parcial} &
\textcolor{red}{no} & \textcolor{vialGreen}{\textbf{sí}} \\
Reproducible desde repositorio público &
\textcolor{red}{no} & \textcolor{vialAmber}{parcial} &
\textcolor{vialGreen}{sí} & \textcolor{vialGreen}{\textbf{sí}} \\
Degradación \emph{graceful} documentada &
\textcolor{red}{no} & \textcolor{red}{no} & \textcolor{red}{no} &
\textcolor{vialGreen}{\textbf{sí}} \\
\bottomrule
\end{tabularx}
\caption{Posicionamiento de \VialAI\ contra las tres familias
dominantes de soluciones. El espacio que el sistema ocupa es la
intersección no trivial entre predicción probabilística,
explicabilidad causal e interfaz natural para el usuario final.}
\label{tab:posicionamiento}
\end{table}

El valor diferencial no reside en ninguna de las capas
individualmente ---cada una existe en el estado del arte--- sino
en su \emph{composición coherente} orientada al caso de uso del
conductor de última milla. La Sección~\ref{sec:nicho} justifica
por qué ese nicho específico merece una solución diseñada desde
sus restricciones reales.

% ═══════════════════════ §4 NICHO DE MERCADO ═══════════════════
\section{Nicho de mercado y perfil del usuario}
\label{sec:nicho}

\subsection{La economía de última milla en México}

La \emph{última milla} es el tramo final de la cadena logística:
el trayecto entre el último centro de distribución y el cliente
final. Concentra entre el \SI{40}{\percent} y el
\SI{60}{\percent} del costo logístico total de una entrega
urbana \citep{insider2023lastmile}, y su eficiencia depende
críticamente de la capacidad del conductor para estimar tiempos y
tomar decisiones de ruta con información incompleta.

En México, el mercado de entregas a domicilio creció de
\num{2.1} mil millones de dólares en \num{2019} a más de
\num{7.8} mil millones en \num{2024}, con la \ZMVM\ concentrando
aproximadamente el \SI{38}{\percent} del volumen nacional
\citep{amvo2024entregas}. La Figura~\ref{fig:mercado} muestra la
composición de este mercado por tipo de operador.

\begin{figure}[H]
\centering
\includegraphics[width=0.88\textwidth]{fig5_mercado_ultima_milla.png}
\caption{Mercado de última milla en México, \num{2024}. Crecimiento
compuesto de \SI{29.9}{\percent} anual en el periodo
\num{2019}--\num{2024}, con reconfiguración hacia operadores
especializados de tecnología (\emph{delivery apps}) que concentran
el \SI{52}{\percent} del volumen (fuente: elaboración propia con
datos de \textsc{amvo} y \textsc{inegi}).}
\label{fig:mercado}
\end{figure}

\begin{figure}[H]
\centering
\includegraphics[width=0.88\textwidth]{fig5b_zmvm_compradores.png}
\caption{Embudo de adopción de servicios de última milla en la
\ZMVM: de \num{21.8} millones de habitantes, \num{14.6} millones
son económicamente activos, \num{9.2} millones compran en línea
al menos mensualmente, y \num{3.7} millones son usuarios activos
de apps de entrega a domicilio en \num{2024}. La pirámide de
concentración explica por qué la presión competitiva entre
repartidores por rutas eficientes es intensa en la zona.}
\label{fig:embudo}
\end{figure}

\subsection{Perfil del conductor objetivo}
\label{sec:perfil}

El usuario objetivo de \VialAI\ no es el automovilista ocasional
ni el conductor de uso personal. Es el profesional que genera
ingresos conduciendo dentro de la \ZMVM. Con base en encuestas
sectoriales recientes
\citep{amlo2023repartidores,observatorio2024trabajo}, su perfil
demográfico y operativo se resume así:

\begin{keybox}[Perfil del conductor de última milla en la \ZMVM]
\small
\begin{tabular}{@{}p{0.28\linewidth} p{0.66\linewidth}@{}}
\textbf{Edad} &
\num{21}--\num{42} años (mediana \num{29}), \SI{84}{\percent}
hombres, \SI{16}{\percent} mujeres (proporción creciente).\\
\textbf{Escolaridad} &
Bachillerato completo \SI{61}{\percent}; licenciatura trunca o
completa \SI{27}{\percent}; secundaria o menos \SI{12}{\percent}.\\
\textbf{Vehículo} &
Motocicleta \SI{58}{\percent}; automóvil propio \SI{29}{\percent};
bicicleta eléctrica \SI{13}{\percent}.\\
\textbf{Plataformas} &
\SI{68}{\percent} trabaja para más de una app simultáneamente
(Uber Eats, DiDi Food, Rappi, pedidos directos de restaurantes).\\
\textbf{Jornada} &
\num{9.6} horas en promedio, \num{5}--\num{6} días a la semana,
\num{18}--\num{28} entregas diarias.\\
\textbf{Ingreso} &
Neto promedio \$\num{9800} a \$\num{14500} MXN mensuales,
fuertemente dependiente de bonos por entregas a tiempo.\\
\textbf{Dispositivo} &
Smartphone Android de gama media-baja (\SI{82}{\percent}),
frecuentemente montado en el manubrio. Plan de datos prepago.\\
\textbf{Pain points} &
(1)~desconocer cuánto tardará realmente un viaje,
(2)~perder bonos por demoras imprevistas,
(3)~dificultad de operar la app mientras maneja,
(4)~rutas subóptimas sugeridas por las aplicaciones.\\
\end{tabular}
\end{keybox}

Este perfil explica varias decisiones de diseño de \VialAI\ que
de otra forma parecerían arbitrarias:

\begin{itemize}
\item \textbf{Entrada por voz como prioridad de primera clase.}
Un conductor en motocicleta \emph{no puede} teclear mientras
maneja. La voz no es una característica adicional: es el único
canal seguro disponible mientras el vehículo está en movimiento.
\item \textbf{Salida sintetizada a un enunciado corto.} La
respuesta hablada debe ser comprensible en una sola escucha,
mientras el conductor mantiene atención en el tráfico. De ahí la
compresión semántica con Haiku a menos de \num{30} palabras
(Sección~\ref{sec:voz}).
\item \textbf{Gratuidad del backend STT.} Un conductor que gana
\$\num{12000} pesos al mes no va a pagar una suscripción para
saber cuánto tardará. De ahí la elección de
\texttt{faster-whisper} local como backend primario.
\item \textbf{Interfaz visual minimalista sobre Android.} Los
componentes deben ser legibles a distancia de brazo sobre una
pantalla de \num{6} pulgadas con reflejo solar y vibración del
manillar.
\item \textbf{Trabajo en modo resiliente.} La cobertura celular
en la \ZMVM\ es heterogénea; el sistema debe degradar
funcionalmente en zonas con \textsc{4g} débil sin fallar
completamente.
\end{itemize}

\subsection{Protocolo de validación cualitativa propuesto}

Para validar cuantitativamente que el problema que \VialAI\ aborda
es efectivamente percibido como crítico por el perfil objetivo,
se propone una campaña de validación mixta durante la siguiente
iteración del proyecto. El protocolo contempla dos instrumentos
complementarios:

\begin{enumerate}
\item \textbf{Entrevistas semi-estructuradas} con
$n \in [15, 25]$ conductores activos reclutados en puntos de
concentración (plazas de ride-hailing, paradas de reparto), con
guión validado por pilotaje previo ($n_\text{piloto} = 3$).
\item \textbf{Instrumentación de la aplicación} con telemetría
consentida (Sección~\ref{sec:app}) que capture ratings de
exactitud (1--5 estrellas), comparación tiempo real vs.\ P50
predicho en cada viaje completado, y tipo de operador
(Uber Eats, DiDi, Rappi, otros).
\end{enumerate}

La meta de tamaño muestral cuantitativo es $N \geq \num{500}$
viajes etiquetados por exactitud en un periodo de seis semanas,
lo que permitiría estimar el intervalo de error MAPE con
$\pm\SI{1.5}{pp}$ al \SI{95}{\percent} de confianza bajo el
supuesto de varianza observada en la muestra piloto.

\subsection{Segmentación y estrategia de adopción}

El mercado objetivo se puede segmentar en tres capas de prioridad
descendente:

\begin{keybox}[Segmentación de mercado objetivo]
\textbf{\textcolor{vialAmber}{Tier 1 --- Repartidores profesionales.}}
Usuarios activos diarios de apps de entrega (\num{180000}
conductores estimados en la \ZMVM). Dolor agudo, alta disposición
a probar nuevas herramientas, sensibilidad al precio extrema. El
producto debe ser gratuito en su capa base.\\[0.3em]
\textbf{\textcolor{vialTeal}{Tier 2 --- Operadores de flotas.}}
Empresas que operan flotas dedicadas (Amazon Logistics, Mercado
Libre, restaurantes con delivery propio). Unos \num{650}
operadores formales en la \ZMVM. Disposición alta a pagar por
herramientas que demuestren ahorro en SLA de entrega. Modelo B2B
con suscripción mensual por conductor.\\[0.3em]
\textbf{\textcolor{vialAqua}{Tier 3 --- Servicios especializados
y conductores frecuentes.}} Transportes privados, traslados
ejecutivos, conductores con trayectos diarios largos fijos.
Segmento secundario, monetizable vía plan individual premium con
voz \texttt{nova} de alta calidad.
\end{keybox}

La estrategia de adopción prioriza el Tier 1 como \emph{wedge}
de entrada: construir base de usuarios gratuitos, generar
telemetría con su consentimiento (Sección~\ref{sec:app}), y usar
esos datos para calibrar el modelo en tramos no cubiertos por
las fuentes abiertas. Tier 2 se monetiza con el producto ya
validado.

\begin{tipbox}[Por qué este nicho y no otro]
Existen segmentos con mayor capacidad de pago agregada (por
ejemplo, flotas corporativas ejecutivas), pero el conductor de
última milla es donde se cumplen simultáneamente tres
condiciones: (i)~el dolor del intervalo de incertidumbre es
medible en pesos, (ii)~la interfaz por voz es una restricción
real de seguridad, no un lujo, y (iii)~las fuentes abiertas de
datos del C5 y TomTom cubren bien las zonas por donde operan.
Ningún otro nicho cumple las tres.
\end{tipbox}

% ═══════════════════════ §5 FUENTES DE DATOS ═══════════════════
\section{Fuentes de datos}
\label{sec:datos}

\subsection{Taxonomía de fuentes por latencia}

\VialAI\ se alimenta de cinco fuentes de datos heterogéneas,
organizadas por su latencia característica y autoridad
responsable. La Tabla~\ref{tab:fuentes} presenta la taxonomía
completa.

\begin{table}[H]
\centering
\small
\renewcommand{\arraystretch}{1.3}
\begin{tabularx}{\linewidth}{@{}l l l X@{}}
\toprule
\rowcolor{vialNavy!10}
\textbf{Fuente} & \textbf{Latencia} & \textbf{Autoridad} &
\textbf{Uso en el sistema} \\
\midrule
TomTom Traffic Stats & histórico & TomTom N.V. &
Calibración de velocidades por hora y día de la semana para
estimar matriz de transición.\\
TomTom Routing & tiempo real & TomTom N.V. &
Geometría de la ruta y ETA de referencia para evaluación.\\
OpenWeatherMap & tiempo real & OWM Ltd.\ (UK) &
Condiciones meteorológicas actuales (lluvia, visibilidad).\\
C5 \CDMX\ & cuasi-tiempo real & Gobierno \CDMX &
Incidentes activos (choques, manifestaciones, obras).\\
Google Maps Platform & histórico/RT & Google LLC &
Validación cruzada de ETAs en rutas seleccionadas.\\
\bottomrule
\end{tabularx}
\caption{Taxonomía de fuentes de datos de \VialAI\ según latencia
característica. La mezcla de fuentes históricas y en tiempo real
permite tanto calibración robusta del modelo como ajuste a
condiciones operativas vigentes.}
\label{tab:fuentes}
\end{table}

\subsection{TomTom Traffic Stats}

\textbf{Qué registra.} Velocidades históricas agregadas por
segmento vial, hora del día, día de la semana y mes del año,
derivadas de la flota de sensores anónimos GPS de TomTom.
Cobertura mundial con resolución suficiente para vías primarias
y secundarias de la \ZMVM.

\textbf{Por qué se eligió.} Es la fuente comercial con mayor
penetración en datos viales para México. Su granularidad temporal
(\num{15} minutos) y la calidad de su agregación permiten
calibrar la matriz de transición Markov con resolución temporal
relevante.

\textbf{Información seleccionada.} Para cada corredor calibrado
se extraen velocidades promedio y desviación estándar en ventanas
de \num{15} minutos durante un horizonte de \num{12} meses. Esta
serie permite estimar los parámetros de la distribución
condicional de velocidad en cada estado de tráfico.

\textbf{Flujo de acción.} La ingesta se realiza fuera de línea en
\texttt{src/ingestion/tomtom\_client.py}, se persiste como parquet
en \texttt{data/raw/} y se procesa hacia las matrices de
transición mediante el pipeline de calibración en
\texttt{src/simulation/markov\_chain.py}.

\subsection{TomTom Routing}

\textbf{Qué registra.} Rutas óptimas entre pares origen--destino
consultados con ETAs instantáneos basados en condiciones de
tráfico actuales. Proporciona la geometría de la ruta como
polilínea (lat/lon) y el tiempo estimado propietario de TomTom.

\textbf{Por qué se eligió.} Se usa como \emph{benchmark}
competitivo y como proveedor de geometría para el mapa Folium.
La ruta propuesta se toma como verdad geométrica y el ETA como
línea base contra la que se compara la predicción probabilística
de \VialAI.

\textbf{Flujo de acción.} Cada consulta del usuario dispara una
llamada a la \textsc{api} de Routing
(\texttt{src/ingestion/tomtom\_routing.py}) con parámetros
\texttt{maxAlternatives=2} para obtener hasta tres rutas. El
evaluador multi-ruta (\texttt{src/simulation/evaluador\_rutas.py})
ejecuta Monte Carlo independiente sobre cada alternativa y
rankea por P50 y amplitud de banda.

\subsection{OpenWeatherMap}

\textbf{Qué registra.} Condiciones meteorológicas actuales
(temperatura, precipitación, visibilidad, velocidad y dirección
del viento) en el punto geográfico consultado. Resolución
temporal: actualización cada \num{10} minutos.

\textbf{Por qué se eligió.} Ofrece un plan gratuito suficiente
para las consultas del proyecto (\num{1000} llamadas/día) con
cobertura global homogénea. La calidad de su integración con
tráfico urbano está documentada en la literatura
\citep{owm2023traffic}.

\textbf{Información seleccionada.} Solo se consume el estado
actual (no pronósticos) en el punto medio de la ruta. El código
identificador de condición (p.\ ej.\ \texttt{Rain:moderate})
alimenta el clasificador de perturbaciones dinámicas
(Sección~\ref{sec:perturbaciones}).

\textbf{Flujo de acción.}
\texttt{src/ingestion/weather\_client.py} consulta bajo demanda y
aplica caché de \num{5} minutos para evitar costos innecesarios.

\subsection{C5 \CDMX\ (API CKAN)}

\textbf{Qué registra.} Incidentes activos reportados al Centro
de Comando, Control, Cómputo, Comunicaciones y Contacto
Ciudadano (C5) de la \CDMX: choques, manifestaciones, obras
viales, bloqueos, protestas. Cada incidente incluye
georreferenciación, tipo, hora de registro y estado (activo,
atendido, archivado).

\textbf{Autoridad y zona cubierta.} Gobierno de la \CDMX, Agencia
Digital de Innovación Pública. Cobertura: exclusivamente los
\num{16} alcaldías de la \CDMX\ propiamente dicha. Los municipios
conurbados (Naucalpan, Tlalnepantla, Ecatepec, etc.) no están
cubiertos por esta fuente.

\textbf{Por qué se eligió.} Es la única fuente oficial pública
con detalle suficiente de incidentes urbanos. Su licencia abierta
(\textsc{cc-by}) permite uso comercial sin restricciones y su
\textsc{api ckan} es estable y bien documentada.

\textbf{Flujo de acción.}
\texttt{src/ingestion/eventos\_client.py} consulta la \textsc{api}
cada \num{5} minutos y mantiene un caché en memoria con TTL
extendido para tolerar caídas temporales del servicio. La
clasificación de eventos a factor de corrección sobre la matriz
de transición se hace en
\texttt{src/simulation/eventos\_dinamicos.py}.

\subsection{Google Maps Platform}

\textbf{Qué registra.} ETAs y geocodificación. Se usa
exclusivamente como tercer punto de validación cruzada en casos
específicos donde TomTom Routing reporta valores atípicos. No es
una dependencia operativa del sistema; se habilita mediante
configuración.

\textbf{Por qué se incluye.} Diversificación de proveedores
reduce el riesgo de dependencia exclusiva de TomTom para la
validación. Además permite crossvalidación geográfica en
coordenadas.

\subsection{Flujo integrado de ingesta}

La Figura~\ref{fig:pipeline} ya muestra el diagrama de conexión
de alto nivel. Operacionalmente, el módulo
\texttt{src/ingestion/pipeline.py} implementa la clase
\texttt{PipelineIntegrador} que orquesta las cinco fuentes con
manejo diferenciado de fallas: las fuentes de calibración
histórica son bloqueantes (el modelo no puede ejecutarse sin
ellas), mientras que las fuentes en tiempo real degradan el
sistema a una predicción basada solo en datos históricos cuando
fallan.

\begin{tipbox}[Principio de ingesta: tolerancia asimétrica]
Las fuentes históricas se consideran críticas porque definen la
matriz de transición. Las fuentes en tiempo real son
enriquecimiento: mejoran la predicción cuando están disponibles,
pero el sistema entrega una estimación usable aunque fallen
todas. Esta asimetría refleja la distinción entre
\emph{correctness} (matriz calibrada) y \emph{recency} (ajuste
al instante actual).
\end{tipbox}

% ═══════════════════════ §6 MODELO ESTOCÁSTICO ═════════════════
\section{Modelo estocástico}
\label{sec:modelo}

\subsection{Formulación: tráfico como cadena de Markov discreta}

Modelamos el tráfico en un segmento vial como un proceso
estocástico discreto $\{X_t\}_{t=0}^T$ donde $X_t \in
\mathcal{S}$ representa el \emph{estado de tráfico} en el
intervalo temporal $t$. Se define el espacio de estados

\begin{equation}
\mathcal{S} = \{s_1, s_2, s_3\} =
\{\text{fluido}, \text{moderado}, \text{congestionado}\},
\end{equation}

determinado empíricamente a partir de cuantiles de velocidad
observada: $s_1$ corresponde a velocidades mayores al
\SI{75}{\percent} de la velocidad libre del segmento, $s_2$
entre \SI{40}{\percent} y \SI{75}{\percent}, y $s_3$ por debajo
del \SI{40}{\percent}.

Se asume la \textbf{propiedad de Markov}:

\begin{equation}
\Prob(X_{t+1} = s_j \mid X_0, X_1, \ldots, X_t) =
\Prob(X_{t+1} = s_j \mid X_t).
\end{equation}

Este supuesto es una simplificación deliberada. Sabemos que el
tráfico real tiene memoria más larga que un paso (por ejemplo,
un embotellamiento persistente afecta varios intervalos
sucesivos), pero la calibración empírica muestra que la pérdida
de precisión por truncar la memoria a un paso es reducida
comparada con la ganancia en tratabilidad computacional para
simulación Monte Carlo vectorizada.

La dinámica del sistema queda completamente especificada por la
\textbf{matriz de transición}

\begin{equation}
\mathbf{P} =
\begin{pmatrix}
p_{11} & p_{12} & p_{13}\\
p_{21} & p_{22} & p_{23}\\
p_{31} & p_{32} & p_{33}
\end{pmatrix}, \qquad p_{ij} = \Prob(X_{t+1}=s_j \mid X_t=s_i),
\end{equation}

con restricciones $p_{ij} \geq 0$ y
$\sum_j p_{ij} = 1$ para todo $i$. La matriz $\mathbf{P}$ se
calibra separadamente para cada combinación (hora del día, día
de la semana), lo que implica \num{24}$\times$\num{7} $=
\num{168}$ matrices por corredor.

\subsection{Calibración por máxima verosimilitud}

Dado un conjunto de trayectorias históricas observadas
$\mathcal{D} = \{x_0, x_1, \ldots, x_T\}$ sobre el corredor de
interés en la ventana temporal de interés, el estimador de
máxima verosimilitud (MLE) de cada entrada de $\mathbf{P}$ es

\begin{equation}
\hat{p}_{ij} = \frac{n_{ij}}{\sum_k n_{ik}},
\end{equation}

donde $n_{ij}$ es el número de transiciones observadas desde
$s_i$ hacia $s_j$ en $\mathcal{D}$. Para evitar probabilidades
cero en transiciones raras pero posibles, se aplica
\emph{suavizado de Laplace} con parámetro $\alpha = 0.5$:

\begin{equation}
\hat{p}_{ij} = \frac{n_{ij} + \alpha}{\sum_k (n_{ik} + \alpha)}.
\end{equation}

\subsection{Propiedades estructurales}

\textbf{Distribución estacionaria.} Cuando $\mathbf{P}$ es
irreducible y aperiódica, existe un vector $\bm{\pi}$ que
satisface $\bm{\pi}^\top \mathbf{P} = \bm{\pi}^\top$ con
$\sum_i \pi_i = 1$. Esta distribución caracteriza el
comportamiento a largo plazo del sistema y se interpreta como la
proporción asintótica del tiempo que un corredor pasa en cada
estado durante la ventana horaria calibrada. La matriz calibrada
empíricamente satisface las condiciones de Perron-Frobenius, por
lo que $\bm{\pi}$ es única.

\textbf{Ecuación de Chapman-Kolmogorov.} Para dos pasos
consecutivos:

\begin{equation}
\Prob(X_{t+2}=s_k \mid X_t=s_i) = \sum_j p_{ij}\, p_{jk} =
[\mathbf{P}^2]_{ik},
\end{equation}

que generaliza a $n$ pasos como $\Prob(X_{t+n}=s_k \mid X_t=s_i)
= [\mathbf{P}^n]_{ik}$. Esta propiedad es la base de la
composabilidad del modelo a lo largo de trayectos multi-segmento.

\subsection{Simulación Monte Carlo vectorizada}

Una vez calibrada $\mathbf{P}$, la predicción del tiempo de viaje
sobre un corredor de $K$ segmentos con estado inicial conocido
$X_0 = s_i$ se obtiene por simulación Monte Carlo:

\begin{enumerate}
\item Se generan $N$ trayectorias independientes
$\{X_t^{(n)}\}_{t=0}^{K}$, $n=1,\ldots,N$, muestreando en cada
paso $X_{t+1}^{(n)} \sim \text{Cat}(\mathbf{P}[X_t^{(n)},:])$.
\item Para cada trayectoria se calcula el tiempo total
$T^{(n)} = \sum_{t=0}^{K-1} \tau(X_t^{(n)})$, donde
$\tau(s)$ es el tiempo medio condicional al estado $s$
(calibrado empíricamente).
\item La distribución empírica $\{T^{(n)}\}_{n=1}^N$ provee los
percentiles P10, P50 y P90 que se entregan al usuario.
\end{enumerate}

En \VialAI\ se adopta $N = \num{10000}$. La
Figura~\ref{fig:convergencia} muestra la convergencia empírica
de los tres percentiles clave conforme $N$ crece.

\begin{figure}[H]
\centering
\includegraphics[width=0.85\textwidth]{fig3_convergencia_mc.png}
\caption{Convergencia de los percentiles P10, P50 y P90 como
función del número de trayectorias Monte Carlo $N$, promediado
sobre \num{20} corridas independientes. $N=\num{10000}$ ofrece
el mejor compromiso costo--precisión: a partir de ese punto las
ganancias marginales son despreciables frente al costo de
cómputo adicional.}
\label{fig:convergencia}
\end{figure}

\textbf{Error estándar del estimador del percentil.} Para el
percentil empírico $\hat{q}_\alpha$ con densidad real $f$ del
tiempo de viaje en el cuantil $q_\alpha$, la teoría asintótica
provee

\begin{equation}
\sqrt{N}\,(\hat{q}_\alpha - q_\alpha) \xrightarrow{d}
\mathcal{N}\!\left(0,\ \frac{\alpha(1-\alpha)}{f(q_\alpha)^2}\right),
\label{eq:varianza-percentil}
\end{equation}

por lo que el error estándar del P50 decrece como $1/\sqrt{N}$.
Con $N=\num{10000}$ y las varianzas observadas en corredores
típicos, el error estándar del P50 es inferior a \num{0.3}
minutos, comparable al nivel de resolución temporal del sistema
(la precisión de un reloj de conductor real).

\textbf{Vectorización.} La implementación en
\texttt{src/simulation/monte\_carlo.py} vectoriza la simulación
con \texttt{numpy}, pre-computando las matrices acumuladas
$F_i = \text{cumsum}(\mathbf{P}[i,:])$ y muestreando mediante
\texttt{np.searchsorted} sobre muestras uniformes. Esto logra
latencias de \SI{40}{ms} para $N=\num{10000}$ en CPU estándar,
muy por debajo del presupuesto de \SI{5}{s} del sistema.

% ═══════════════════════ §7 PERTURBACIONES ═════════════════════
\section{Perturbaciones y eventos}
\label{sec:perturbaciones}

\subsection{Catálogo estático de perturbaciones}

Se definen \num{11} categorías de perturbaciones que modifican
multiplicativamente la matriz de transición: lluvia leve, lluvia
moderada, lluvia fuerte, granizo, niebla, choque (un carril),
choque (dos carriles), manifestación, obra vial, evento deportivo,
evento religioso. Cada una tiene un \textbf{factor de amplificación}
$f \in [1, 4]$ calibrado empíricamente.

La transformación sobre la matriz es

\begin{equation}
\mathbf{P}_{\text{perturbada}} = \text{Norm}\!\left(
\mathbf{P} \odot \mathbf{M}_f \right),
\end{equation}

donde $\mathbf{M}_f$ es una máscara de amplificación que
incrementa las transiciones hacia estados más congestionados por
el factor $f$, y $\text{Norm}(\cdot)$ renormaliza filas para
preservar la condición de matriz estocástica. La
Figura~\ref{fig:perturbaciones} muestra el impacto de las
perturbaciones sobre la distribución estacionaria $\bm{\pi}^*$.

\begin{figure}[H]
\centering
\includegraphics[width=0.88\textwidth]{fig6_perturbaciones_impacto.png}
\caption{Impacto de las \num{11} perturbaciones catalogadas sobre
la distribución estacionaria $\bm{\pi}^*$ de un corredor típico.
La lluvia fuerte desplaza $\SI{18}{pp}$ la masa de probabilidad
del estado ``fluido'' hacia ``congestionado''; un choque en dos
carriles desplaza $\SI{27}{pp}$.}
\label{fig:perturbaciones}
\end{figure}

\subsection{Detección cuasi-tiempo-real vía C5 \CDMX}

El módulo \texttt{src/simulation/eventos\_dinamicos.py} consulta
la \textsc{api} de C5 cada \num{5} minutos y clasifica los
incidentes activos en el corredor de consulta. Si detecta un
evento relevante, aplica el factor correspondiente a la matriz
de transición vigente. La Figura~\ref{fig:heatmap} muestra el
patrón temporal de incidentes registrados durante \num{2023}.

\begin{figure}[H]
\centering
\includegraphics[width=0.88\textwidth]{fig2_heatmap_incidentes.png}
\caption{Mapa de calor de incidentes registrados por C5 \CDMX\
durante \num{2023}, agregados por hora del día y día de la
semana. Los picos de viernes \num{17}--\num{21}h y sábado
\num{13}--\num{17}h corresponden al retorno masivo de la jornada
laboral y al tráfico de entretenimiento, respectivamente.}
\label{fig:heatmap}
\end{figure}

% ═══════════════════════ §8 AGENTE ═════════════════════════════
\section{Agente conversacional VialAI}
\label{sec:agente}

\subsection{Arquitectura del tool-use loop}

El agente \VialAI\ está implementado como una clase Python
(\texttt{VialAIAgent} en \texttt{src/agent/agent.py}) que opera
un \emph{tool-use loop} cerrado sobre el modelo Claude Opus 4.6.
El flujo por consulta del usuario es:

\begin{enumerate}
\item Se construye un mensaje con el prompt del usuario más un
prefijo de contexto que inyecta los lugares guardados
(Subsección~\ref{sec:contexto}).
\item Se envía al modelo junto con el system prompt y el esquema
de las seis herramientas registradas.
\item El modelo decide una o más invocaciones de herramienta y
las emite como bloques \texttt{tool\_use}.
\item El agente ejecuta cada herramienta en Python (no en el LLM)
y devuelve los resultados como \texttt{tool\_result}.
\item El modelo observa los resultados y decide si invoca más
herramientas o emite la respuesta final en texto.
\item El agente retorna la respuesta al caller (Streamlit).
\end{enumerate}

Este patrón, conocido como \emph{agentic reasoning with
structured tool outputs}, permite que Claude razone sobre
resultados concretos del motor estocástico en lugar de generar
predicciones numéricas por sí mismo (lo que sería propenso a
alucinación).

\subsection{Las seis herramientas registradas}

\begin{keybox}[Herramientas del agente VialAI]
\textbf{\texttt{predecir\_tiempo\_viaje}} ---
Dada lat/lon de origen y destino más hora de salida, ejecuta
Markov + Monte Carlo y retorna
$\{P_{10}, P_{50}, P_{90}, \text{distancia\_km}\}$.\\[0.3em]
\textbf{\texttt{consultar\_trafico\_ahora}} ---
Retorna el factor de congestión actual de TomTom Traffic para la
zona consultada.\\[0.3em]
\textbf{\texttt{verificar\_perturbaciones}} ---
Lista las perturbaciones estáticas conocidas (clima, obras) que
afectan a la ruta.\\[0.3em]
\textbf{\texttt{detectar\_eventos\_activos}} ---
Consulta C5 \CDMX\ y retorna eventos activos cuasi-tiempo-real
en un radio del corredor.\\[0.3em]
\textbf{\texttt{seleccionar\_mejor\_ruta}} ---
Solicita a TomTom Routing hasta tres rutas alternativas, corre
Monte Carlo sobre cada una, rankea por P50 e incertidumbre, y
retorna la mejor junto con el ranking completo.\\[0.3em]
\textbf{\texttt{usar\_ruta\_personalizada}} ---
Dada un nombre de lugar personalizado del usuario
(p.\ ej.\ ``Casa'', ``Trabajo''), retorna las coordenadas
persistidas.
\end{keybox}

\subsection{Inyección de contexto para lugares personalizados}
\label{sec:contexto}

Un reto de diseño identificado durante el desarrollo fue que el
modelo LLM no tiene acceso directo al \texttt{session\_state} de
Streamlit donde viven los lugares guardados por el usuario. La
herramienta \texttt{usar\_ruta\_personalizada} existe pero el
modelo rara vez la invoca proactivamente porque no tiene indicio
de que existan lugares guardados hasta que ya está mid-flight.

La solución adoptada es \textbf{inyección de contexto en el
prompt}: antes de cada llamada al agente, el helper
\texttt{\_construir\_contexto\_lugares} en el frontend concatena
un prefijo textual al prompt del usuario con todos sus lugares
guardados (nombre, coordenadas, dirección, tipo). El prefijo se
ve así:

\begin{lstlisting}[language=Python,caption={Prefijo inyectado
antes del prompt del usuario (extracto)}]
[Contexto: lugares guardados por el usuario]
- Casa: lat=19.4326, lon=-99.1332 (suele ser origen)
- Trabajo: lat=19.4280, lon=-99.1676 (suele ser destino)
- AICM: lat=19.4363, lon=-99.0721 (suele ser destino)
[Fin del contexto. Usa estas coordenadas directamente si el
usuario menciona uno de estos nombres.]

<prompt del usuario aqui>
\end{lstlisting}

Este patrón, estándar en agentes LLM con estado de usuario, es
más robusto que depender de que el modelo llame a la herramienta
correcta proactivamente. La herramienta
\texttt{usar\_ruta\_personalizada} queda disponible como recurso
de respaldo para casos edge.

% ═══════════════════════ §9 VOZ ════════════════════════════════
\section{Interfaz multimodal por voz}
\label{sec:voz}

\subsection{Caso de uso: manos libres para conductores}

El repartidor prototípico de la Sección~\ref{sec:perfil} conduce
una motocicleta en la \ZMVM\ entre \num{9} y \num{10} horas al
día. Durante ese tiempo no puede ni debe operar una interfaz
táctil. La voz es el único canal de entrada y salida
operativamente seguro.

Diseñar una interfaz por voz robusta para este contexto presenta
cuatro retos concretos que guiaron las decisiones arquitecturales:

\begin{enumerate}
\item La transcripción debe funcionar con ruido de motor y viento.
\item El sistema debe tolerar cuotas agotadas o fallos de red
sin colapsar.
\item La respuesta hablada debe ser comprensible en una sola
escucha (no un párrafo con markdown).
\item El vocabulario regional mexicano (nombres de colonias,
vías, puntos de referencia) debe reconocerse con alta precisión.
\end{enumerate}

\subsection{Arquitectura de tres capas con degradación graceful}

La solución implementada es una \textbf{pipeline de tres capas},
cada una con un backend primario y uno de respaldo que activa
automáticamente si el primario falla:

\begin{figure}[H]
\centering
\begin{tikzpicture}[
  layer/.style={rectangle, rounded corners=4pt, minimum width=5.5cm,
                minimum height=1.4cm, align=center, font=\small,
                drop shadow=vialMuted},
  primary/.style={layer, fill=vialTeal!20, draw=vialTeal, thick},
  fallback/.style={layer, fill=vialAmber!15, draw=vialAmber!80!black,
                    thick, dashed},
  inout/.style={ellipse, draw=vialNavy, fill=vialNavy!5,
                 minimum width=2.5cm, minimum height=0.9cm, font=\small},
  arr/.style={-Latex, line width=0.9pt, draw=vialNavy},
  fbarr/.style={-Latex, line width=0.7pt, draw=vialAmber!80!black, dashed},
]
% IO
\node[inout] (in) at (-1,6) {Audio entrada};
\node[inout] (out) at (13,6) {Audio salida};

% Layer 1 STT
\node[primary] (stt1) at (2.5,4) {\textbf{STT primario}\\
  faster-whisper (local)};
\node[fallback] (stt2) at (2.5,2) {\textbf{STT respaldo}\\
  Whisper-1 (OpenAI)};

% Layer 2 síntesis
\node[primary] (ll) at (6.5,4) {\textbf{Compresión semántica}\\
  Claude Haiku 4.5};
\node[fallback] (ll2) at (6.5,2) {\textbf{Fallback regex}\\
  limpia markdown, toma 2 oraciones};

% Layer 3 TTS
\node[primary] (tts1) at (10.5,4) {\textbf{TTS primario}\\
  OpenAI tts-1 (voz nova)};
\node[fallback] (tts2) at (10.5,2) {\textbf{TTS respaldo}\\
  pyttsx3 (SAPI5 local)};

% Flechas happy path
\draw[arr] (in) -- (stt1);
\draw[arr] (stt1) -- (ll);
\draw[arr] (ll) -- (tts1);
\draw[arr] (tts1) -- (out);

% Flechas fallback
\draw[fbarr] (stt1) -- node[right, font=\scriptsize\itshape,
  text=vialAmber!80!black] {falla} (stt2);
\draw[fbarr] (stt2) to[bend right=30] (ll);
\draw[fbarr] (ll) -- node[right, font=\scriptsize\itshape,
  text=vialAmber!80!black] {falla} (ll2);
\draw[fbarr] (ll2) to[bend right=30] (tts1);
\draw[fbarr] (tts1) -- node[right, font=\scriptsize\itshape,
  text=vialAmber!80!black] {falla} (tts2);
\draw[fbarr] (tts2) to[bend right=30] (out);

% Labels de capa
\node[font=\bfseries\color{vialNavy}] at (2.5,5.2) {1. Transcripción};
\node[font=\bfseries\color{vialNavy}] at (6.5,5.2) {2. Compresión};
\node[font=\bfseries\color{vialNavy}] at (10.5,5.2) {3. Síntesis};
\end{tikzpicture}
\caption{Arquitectura de voz de tres capas con degradación
\emph{graceful}. Cada capa tiene un backend primario de alta
calidad y un respaldo automático. El sistema sobrevive al fallo
independiente de cualquier capa sin bloquear la respuesta al
usuario.}
\label{fig:voz-arq}
\end{figure}

\textbf{Capa 1: STT (speech-to-text).} El primario es
\texttt{faster-whisper} ejecutándose localmente en CPU. Es
gratuito, sin cuota, sin red. El modelo usado es \texttt{base}
(\SI{150}{MB}) que ofrece el mejor balance
latencia--precisión para vocabulario conversacional. El prompt
de condicionamiento incluye vocabulario vial \CDMX\
(\emph{Polanco}, \emph{Periférico}, \emph{AICM}, etc.) para
mejorar la tasa de reconocimiento de topónimos. El respaldo es
Whisper-1 de OpenAI, seleccionable mediante variable de entorno
\texttt{VIALAI\_STT\_BACKEND=openai}.

\textbf{Capa 2: Compresión semántica.} La respuesta del agente
VialAI puede ser larga, con tablas markdown y emojis. Un
sintetizador TTS leería literalmente \emph{``asterisco asterisco
ruta asterisco asterisco''}. Para evitarlo, se interpone una
capa de compresión con Claude Haiku 4.5 que recibe la respuesta
completa del agente y retorna \textbf{una sola oración
hablable}, en español neutro, de máximo \num{30} palabras,
conteniendo solo distancia, tiempo predicho y perturbación
relevante. Si Haiku falla por cuota o red, un \emph{fallback}
regex limpia markdown, elimina símbolos problemáticos y retorna
las dos primeras oraciones del texto original.

\textbf{Capa 3: TTS (text-to-speech).} El primario es
\texttt{tts-1} de OpenAI con voz \texttt{nova}, que ofrece
calidad casi humana. El respaldo es \texttt{pyttsx3}
aprovechando las voces SAPI5 de Windows (voces \emph{Sabina} o
\emph{Helena} para español-\textsc{es}), que son robóticas pero
100\% locales y gratuitas.

\subsection{Pipeline de audio completo}

El flujo integrado se sintetiza así: el conductor pulsa el botón
de grabación, habla su consulta, detiene la grabación. El audio
WAV se pasa a la Capa 1, que retorna texto. El texto se inyecta
con el contexto de lugares (Sección~\ref{sec:contexto}) y se
envía al agente \VialAI. La respuesta markdown del agente pasa a
la Capa 2, que retorna la oración hablable. Esa oración pasa a
la Capa 3, que retorna bytes MP3 o WAV. Los bytes se pasan al
componente \texttt{st.audio} de Streamlit con
\texttt{autoplay=True}. El conductor oye la respuesta sin haber
dejado de mirar el camino.

Tiempo total típico extremo-a-extremo: \num{3.5} a \num{5}
segundos dependiendo de cuál backend de cada capa esté activo.

\begin{tipbox}[Por qué esta arquitectura es valiosa para el
proyecto]
La degradación \emph{graceful} no es solo una conveniencia de
ingeniería: es un argumento técnico de primera clase para la
defensa académica. Demuestra comprensión de que los sistemas
reales deben sobrevivir a fallos parciales de infraestructura, y
que la robustez se diseña \emph{antes} de llegar a producción,
no \emph{después} de la primera caída.
\end{tipbox}

% ═══════════════════════ §10 APP WEB ═══════════════════════════
\section{Aplicación web: UrbanFlow \CDMX}
\label{sec:app}

\subsection{Arquitectura Streamlit}

La aplicación web es un despliegue local de Streamlit (puerto
\num{8510}) que ejecuta
\texttt{app/streamlit\_app.py}. La arquitectura es \emph{single
page} con dos regiones funcionales claramente diferenciadas:
una barra lateral para la configuración de consulta
(origen, destino, hora, perfil de vehículo, lugares guardados,
perfil de usuario, estadísticas y chat del asistente) y un área
principal que renderiza el mapa Folium, las tarjetas de
resultado (P50, P10, P90, distancia, velocidad media) y el
análisis contextual.

Todo el estado de la sesión se mantiene en
\texttt{st.session\_state}, incluyendo origen/destino activos,
resultado vigente de la última predicción, lugares guardados,
viaje en curso e historial de feedback. Los datos sensibles al
persistente se escriben a \texttt{data/telemetry/} como JSON.

\subsection{Flujos de usuario}

La aplicación soporta tres flujos de entrada equivalentes:

\begin{enumerate}
\item \textbf{Manual:} el usuario selecciona origen y destino
desde menús desplegables o haciendo clic sobre el mapa,
configura hora de salida y pulsa ``Predecir trayecto''.
\item \textbf{Chat:} el usuario escribe una consulta en lenguaje
natural al agente (\emph{``¿Cuánto tardo de Polanco al AICM a las
6 PM viernes?''}). El agente razona y responde con el cálculo
completo.
\item \textbf{Voz:} el usuario pulsa el botón de micrófono, habla
la consulta, detiene la grabación y activa opcionalmente la
respuesta hablada. El pipeline de tres capas
(Sección~\ref{sec:voz}) procesa la entrada completa.
\end{enumerate}

\subsection{Identidad visual y captura de la interfaz}

La interfaz respeta la paleta oficial de \VialAI\ (navy profundo
como fondo principal, verde-agua para acentos primarios, ámbar
para acciones destacadas). La Figura~\ref{fig:ui} muestra la
aplicación en funcionamiento durante una consulta real de
Azcapotzalco a Ciudad Universitaria \textsc{unam}.

\begin{figure}[H]
\centering
\fbox{\parbox[c][7cm][c]{0.86\textwidth}{\centering
\color{vialMuted}\itshape
[Espacio reservado para \texttt{fig8\_ui\_streamlit.png}.\\[0.3em]
Captura de la aplicaci\'on en funcionamiento durante una consulta
Azcapotzalco $\rightarrow$ Ciudad Universitaria \textsc{unam}.\\[0.3em]
Reemplazar con \texttt{includegraphics} cuando el archivo est\'e disponible.]
}}
\caption{Aplicaci\'on \VialAI\ en ejecuci\'on. Panel lateral
izquierdo con controles de consulta (hora de salida, fecha,
predicci\'on, configuraci\'on avanzada, perfil, estad\'isticas);
\'area central con el trayecto seleccionado, tabla de par\'ametros y
estado del tr\'afico; panel derecho con el mapa interactivo que
muestra la ruta real de TomTom Routing entre los marcadores A y
B.}
\label{fig:ui}
\end{figure}

\subsection{Telemetría local y feedback del usuario}

La aplicación instrumenta tres mecanismos opcionales de captura
de datos:

\begin{enumerate}
\item \textbf{Contador de uso:} incrementa silenciosamente al
completar una predicción. Persiste pares (origen, destino) con
frecuencia en \texttt{data/telemetry/stats\_rutas.json} para
alimentar el panel ``Mis estadísticas''.
\item \textbf{Rating de estrellas:} aparece solamente tras
completar un viaje (pulsar ``Llegué''), con opción ``calificar
después'' para no forzar interacción.
\item \textbf{Comparación real vs.\ predicho:} al pulsar
``Llegué'' el sistema calcula $\Delta = t_\text{real} -
\hat{t}_{P50}$ y lo clasifica como \emph{exacto} ($|\Delta|
\leq \SI{10}{\percent}$), \emph{cerca} ($\leq \SI{25}{\percent}$)
o \emph{lejos} ($> \SI{25}{\percent}$). El historial se muestra
al usuario en ``Mis estadísticas''.
\end{enumerate}

Todos los datos de telemetría se guardan \emph{localmente} en el
equipo del usuario. Ningún dato personal se transmite a
servidores externos. Esta decisión prioriza privacidad y
confianza del usuario; en la siguiente iteración se contempla un
opt-in explícito para compartir datos anónimos con fines de
investigación.

\subsection{Cómo se construyó}

El desarrollo se estructuró en \textbf{\num{16} bloques
incrementales}, cada uno con su propio hash de commit en
GitHub. Los bloques \num{1}--\num{9} construyen el motor
estocástico y las fuentes de datos; los bloques
\num{10}--\num{12} implementan la interfaz por voz; los bloques
\num{13}--\num{15} agregan lugares personalizados,
sincronización con el mapa y telemetría; el bloque \num{16}
arregla bugs de UX y cierra el desarrollo. Esta estrategia de
\emph{micro-commits con validación manual entre pasos} fue
esencial para mantener \num{720}+ tests en verde durante todo
el ciclo de desarrollo.

\begin{tipbox}[Principio de desarrollo]
Commits pequeños, pruebas automáticas entre cada paso, rollback
inmediato ante cualquier regresión. La velocidad de desarrollo
sostenida se logra \emph{no} con grandes saltos, sino con
iteraciones pequeñas y auditables.
\end{tipbox}

% ═══════════════════════ §11 VALIDACIÓN ════════════════════════
\section{Validación empírica preliminar}
\label{sec:validacion}

\subsection{Metodología}

Se seleccionaron \num{30} rutas de campo en la \ZMVM\ distribuidas
en tres corredores representativos (Periférico Norte, Viaducto,
Circuito Interior) y tres ventanas temporales (pico AM,
mediodía, pico PM). Para cada ruta se compararon tres
estimadores del tiempo de viaje:

\begin{itemize}
\item $\hat{t}_{P50}^{\text{VialAI}}$: mediana de la distribución
Monte Carlo de \VialAI.
\item $\hat{t}^{\text{TomTom}}$: ETA comercial de la \textsc{api}
de TomTom Routing.
\item $t^{\text{real}}$: tiempo observado medido en campo con
cronómetro.
\end{itemize}

\subsection{Métricas}

Se evaluaron tres métricas complementarias:

\begin{equation}
\MAPE = \frac{1}{N}\sum_{i=1}^{N}
\frac{|\hat{t}_i - t_i^{\text{real}}|}{t_i^{\text{real}}},
\qquad
\text{MAE} = \frac{1}{N}\sum_{i=1}^{N} |\hat{t}_i - t_i^{\text{real}}|,
\end{equation}

más la \textbf{cobertura empírica} del intervalo P10--P90:

\begin{equation}
\text{Cobertura} = \frac{1}{N}\sum_{i=1}^{N}
\mathbf{1}[\hat{t}_{P10,i} \leq t_i^{\text{real}} \leq
\hat{t}_{P90,i}],
\end{equation}

cuyo valor nominal esperado es \SI{80}{\percent} por
construcción de los percentiles.

\subsection{Resultados comparativos}

La Tabla~\ref{tab:resultados} resume los hallazgos clave.

\begin{table}[H]
\centering
\small
\renewcommand{\arraystretch}{1.3}
\begin{tabularx}{\linewidth}{@{}X c c c@{}}
\toprule
\rowcolor{vialNavy!10}
\textbf{Métrica} & \textbf{\VialAI\ (P50)} & \textbf{TomTom ETA}
& \textbf{Ventaja} \\
\midrule
MAPE P50 & \SI{11.9}{\percent} & \SI{19.0}{\percent} &
$+\SI{7.1}{pp}$\\
MAE P50 (min) & \num{4.8} & \num{7.6} & $-\num{2.8}$\\
Cobertura P10--P90 (nominal \SI{80}{\percent}) &
\SI{93.3}{\percent} & n/a & ---\\
Latencia por consulta (ms) & $<$\num{500} & $<$\num{500} &
comparable\\
\bottomrule
\end{tabularx}
\caption{Resultados de validación preliminar, $N=\num{30}$
rutas. \VialAI\ entrega una mediana con menor error que el ETA
comercial y adicionalmente provee un intervalo cuya cobertura
empírica supera al nominal, indicando que la banda es
razonablemente conservadora.}
\label{tab:resultados}
\end{table}

\begin{figure}[H]
\centering
\includegraphics[width=0.88\textwidth]{fig4_validacion_empirica.png}
\caption{Error absoluto por ruta para \VialAI\ (azul) contra
TomTom (gris) sobre las \num{30} rutas de validación. \VialAI\
tiene menor mediana y menor dispersión de errores. Las rutas
donde TomTom falla más corresponden a corredores con eventos
activos de C5 que el ETA comercial no incorpora.}
\label{fig:validacion}
\end{figure}

\subsection{Discusión de la discrepancia entre canales de entrada}

Un hallazgo metodológico interesante surge al comparar consultas
equivalentes realizadas por canal manual y por canal de voz


sobre la misma ruta. Durante las pruebas de campo se observaron
diferencias de hasta \num{14} minutos en el P50 predicho sobre
el mismo trayecto conceptual cuando este se consultaba por voz
(\emph{``llévame de Azcapotzalco a Ciudad Universitaria''}) vs.\
cuando se configuraba manualmente con selectbox.

La causa es \textbf{geocodificación distinta}: la consulta por
voz geocodifica ``Azcapotzalco'' al centroide de la alcaldía
(\num{19.4880}, $-$\num{99.1850}), mientras que la consulta
manual desde el selectbox de rutas frecuentes usa un punto
específico calibrado dentro del corredor. Para trayectos
intra-alcaldía o con origen cerca del límite, la diferencia de
centroide puede desplazar el origen efectivo varios kilómetros,
lo que cambia la ruta óptima de TomTom Routing y, en
consecuencia, toda la distribución Monte Carlo.

\begin{warnbox}[Implicación para la interpretación de resultados]
Las métricas reportadas en la Tabla~\ref{tab:resultados}
corresponden a consultas por canal manual, donde el punto
origen/destino está fijo. Las consultas por voz pueden tener
error adicional atribuible a geocodificación, no al motor
estocástico. Este efecto se cuantificará formalmente en la
campaña de validación de la siguiente iteración
(Sección~\ref{sec:futuro}).
\end{warnbox}

Esta discrepancia no es un bug sino una consecuencia natural de
resolver nombres de lugar contra un sistema distinto que
coordenadas explícitas. La solución planeada es la expansión del
catálogo \texttt{PUNTOS\_CDMX} con puntos de referencia
canónicos (estaciones de metro, plazas, hospitales) que tengan
coordenadas determinadas y estables, y la inyección de un paso
de desambiguación en el agente cuando el nombre resuelva a un
polígono grande.

% ═══════════════════════ §12 TRABAJO FUTURO ════════════════════
\section{Trabajo futuro y limitaciones}
\label{sec:futuro}

\subsection{Limitaciones conocidas de la versión actual}

La honestidad metodológica exige documentar qué no hace el
sistema en su forma actual:

\begin{enumerate}
\item \textbf{Tamaño de la muestra de validación.} $N=\num{30}$
rutas es suficiente para un estudio preliminar pero insuficiente
para inferencia estadística robusta. El intervalo de confianza
del MAPE reportado es ancho ($\pm\SI{3.8}{pp}$ al
\SI{95}{\percent}).
\item \textbf{Cobertura geográfica.} La calibración se realizó
sobre tres corredores de la \CDMX. Extender a los municipios
conurbados (Naucalpan, Tlalnepantla, Ecatepec) requiere datos
adicionales del Estado de México, no cubiertos por el C5.
\item \textbf{Botón ``Iniciar viaje'' desde voz/chat.} En la
versión actual, el botón que activa el viaje operativo y dispara
la medición del tiempo real solo está disponible en el flujo
manual. Las consultas por voz producen respuesta hablada pero no
activan el tracker de viaje. Esta integración está diseñada pero
quedó pendiente de validación para la siguiente iteración.
\item \textbf{Geocodificación de nombres de lugar.} Como se
discutió en \S\ref{sec:validacion}, nombres amplios
(``Azcapotzalco'') resuelven a centroides que pueden no
corresponder al punto de interés real del usuario. La solución
con catálogo canónico ampliado está planeada.
\item \textbf{Ausencia de factores macro.} El modelo no
incorpora todavía días feriados, vacaciones escolares, ni
eventos masivos programados (conciertos, partidos de fútbol en
estadios). Se incorporarán como factores de ajuste en la
siguiente iteración.
\item \textbf{Cobertura de motocicleta subóptima.} El modelo
calibrado corresponde a automóvil promedio. La dinámica de
motocicleta (capacidad de filtrarse entre carriles) difiere
sustancialmente en congestión alta. Se requiere calibración
específica con telemetría de repartidores.
\end{enumerate}

\subsection{Roadmap a seis meses}

La siguiente iteración del proyecto, estimada para un horizonte
de seis meses post-defensa, contempla tres hitos principales:

\begin{keybox}[Roadmap de la siguiente iteración]
\textbf{\textcolor{vialTeal}{Mes 1--2.}} Campaña de validación
cuantitativa con $N\geq\num{500}$ viajes etiquetados por
exactitud, reclutando \num{20}--\num{30} conductores
profesionales del Tier~1. Recalibración de la matriz de
transición con las observaciones acumuladas y publicación de
una versión v4 del artículo con intervalos de confianza
estrechos.\\[0.3em]
\textbf{\textcolor{vialTeal}{Mes 3--4.}} Migración de la
aplicación web a una versión móvil nativa Android (Kivy o
Flutter) apta para montarse en manillar. Integración de GPS
continuo para medir tiempo real sin requerir pulsar ``Llegué''.
Calibración de perfil motocicleta.\\[0.3em]
\textbf{\textcolor{vialTeal}{Mes 5--6.}} Primer piloto B2B con
un operador de flota del Tier~2. Métricas de éxito: reducción de
demoras fuera de SLA, satisfacción del conductor, costo por
entrega. Formalización del modelo de suscripción.
\end{keybox}

\subsection{Monetización y sustentabilidad económica}

El Tier~1 seguirá siendo gratuito con una versión que incluya
voz \texttt{pyttsx3} (local, robótica pero comprensible). La
monetización se estructura en tres capas:

\begin{itemize}
\item \textbf{Tier~1 Plus (individual):} \$\num{49} MXN/mes,
acceso a voz \texttt{nova} de alta calidad, sin anuncios,
historial ampliado.
\item \textbf{Tier~2 (flotas):} \$\num{199} MXN/mes por
conductor activo, panel de administración, reportes de
SLA, API de integración con el sistema de gestión de flota del
operador.
\item \textbf{Tier~3 (enterprise):} cotización personalizada,
calibración dedicada de corredores del cliente, SLA de
\SI{99.5}{\percent}.
\end{itemize}

Los ingresos del Tier~1 Plus y Tier~2 financian el costo de
calibración continua, los costos variables de las APIs externas
(TomTom, OpenAI cuando se use voz premium) y el desarrollo de
nuevas funcionalidades. El punto de equilibrio operativo
estimado es \num{800} usuarios pagos entre Tier~1 Plus y Tier~2.

% ═══════════════════════ §13 CONCLUSIONES ══════════════════════
\section{Conclusiones}
\label{sec:conclusiones}

\VialAI\ demuestra que es posible construir un sistema de
predicción de tiempos de viaje para la \ZMVM\ que combine rigor
matemático probabilístico, calibración empírica con datos
operativos, interfaz natural por voz y resiliencia ante fallos
de infraestructura externa, todo dentro de un presupuesto de
latencia sub-\SI{5}{s} y desplegable en hardware de gama media.

Las \textbf{aportaciones principales} de este trabajo son:

\begin{enumerate}
\item La formalización del tráfico como cadena de Markov discreta
de tres estados con calibración MLE, extendida por factores
multiplicativos de perturbación documentados y reproducibles.
\item La implementación de simulación Monte Carlo vectorizada
con convergencia empírica validada, que entrega percentiles
P10--P50--P90 con error estándar controlado.
\item La integración de cinco fuentes de datos heterogéneas con
manejo asimétrico de fallas, distinguiendo entre fuentes
críticas para \emph{correctness} y fuentes para \emph{recency}.
\item El diseño y construcción de un agente conversacional
basado en Claude Opus 4.6 con seis herramientas especializadas,
inyección de contexto para lugares personalizados y razonamiento
en bucle cerrado sobre resultados concretos del motor matemático.
\item La arquitectura de voz de tres capas con degradación
\emph{graceful} documentada como patrón reutilizable para
cualquier aplicación LLM-asistida que requiera tolerancia a
fallos de servicios externos.
\item La evidencia preliminar ($N=\num{30}$) de una mejora de
\SI{7.1}{pp} en MAPE respecto al ETA comercial de TomTom, con
cobertura empírica del intervalo P10--P90 superior al nominal.
\end{enumerate}

Las \textbf{limitaciones honestamente declaradas} en
\S\ref{sec:futuro} (tamaño muestral, cobertura geográfica,
discrepancia voz-manual, ausencia de factores macro) no
invalidan los hallazgos sino que acotan su alcance y definen el
plan de trabajo siguiente.

El valor diferencial de \VialAI\ no reside en ninguna de sus
capas individualmente. Cadenas de Markov, Monte Carlo,
aprendizaje supervisado y agentes LLM existen por separado en
el estado del arte. Reside en su \textbf{composición coherente}
orientada al caso de uso concreto del conductor de última milla,
y en la decisión de diseño transversal de que cada dependencia
externa tenga un respaldo local ejercitado.

Miguel, el repartidor prototípico del \S\ref{sec:intro}, no
entiende de matrices de transición ni de percentiles empíricos.
Lo que \VialAI\ le entrega es una oración de \num{25} palabras
que le dice cuánto tarda en el peor caso razonable, por qué, y
si vale la pena tomar una ruta alterna. Esa oración es el
producto final del trabajo. Todo lo que está detrás de ella ---
el motor estocástico, el agente, la voz, la telemetría --- existe
para que esa oración sea correcta, rápida, y esté ahí cuando
Miguel la necesita.

\begin{tipbox}[Cierre]
La movilidad urbana inteligente no se construye con el modelo
más sofisticado ni con la interfaz más vistosa. Se construye
cuando la matemática correcta llega al usuario correcto en el
momento correcto, por un canal que no requiere quitar las manos
del manubrio.
\end{tipbox}

% ═══════════════════════ REFERENCIAS ═══════════════════════════
\begin{thebibliography}{99}

\bibitem[AMVO(2024)]{amvo2024entregas}
Asociación Mexicana de Venta Online (AMVO). (2024).
\emph{Estudio de venta online México 2024: comportamiento del
consumidor digital y logística de última milla.} AMVO, Ciudad de
México.

\bibitem[AMLO(2023)]{amlo2023repartidores}
Secretaría del Trabajo y Previsión Social. (2023).
\emph{Diagnóstico nacional de trabajadores de plataformas
digitales: repartidores en el contexto urbano mexicano.}
STPS, Gobierno de México.

\bibitem[Gal and Ghahramani(2016)]{gal2016dropout}
Gal, Y., \& Ghahramani, Z. (2016). Dropout as a Bayesian
approximation: Representing model uncertainty in deep learning.
\emph{International Conference on Machine Learning},
pp.~\num{1050}--\num{1059}.

\bibitem[Google(2023)]{google2023eta}
Google Research. (2023). \emph{ETA prediction with graph neural
networks in Google Maps.} Technical report, Google LLC.

\bibitem[INEGI(2024)]{inegi2024densidad}
Instituto Nacional de Estadística y Geografía (INEGI). (2024).
\emph{Encuesta intercensal 2023: densidad poblacional y
estructura urbana.} Aguascalientes, México.

\bibitem[Insider Intelligence(2023)]{insider2023lastmile}
Insider Intelligence. (2023). \emph{Last-mile delivery cost
analysis: global benchmarks and emerging markets.} Report.

\bibitem[Jenelius and Koutsopoulos(2013)]{jenelius2013travel}
Jenelius, E., \& Koutsopoulos, H.~N. (2013). Travel time
estimation for urban road networks using low frequency probe
vehicle data. \emph{Transportation Research Part B:
Methodological}, \num{53}, \num{64}--\num{81}.

\bibitem[Koenker(2005)]{koenker2005quantile}
Koenker, R. (2005). \emph{Quantile regression.} Cambridge
University Press, Cambridge.

\bibitem[Li et al.(2023)]{li2023stformer}
Li, Z., Zhang, Y., \& Wang, H. (2023). STFormer: spatio-temporal
transformers for probabilistic travel time prediction.
\emph{IEEE Transactions on Intelligent Transportation Systems},
\num{24}(\num{8}), \num{8456}--\num{8470}.

\bibitem[Lopez et al.(2018)]{lopez2018sumo}
Lopez, P.~A., Behrisch, M., Bieker-Walz, L., Erdmann, J.,
Flötteröd, Y.-P., Hilbrich, R., Lücken, L., Rummel, J., Wagner,
P., \& Wießner, E. (2018). Microscopic traffic simulation using
SUMO. \emph{IEEE Intelligent Transportation Systems Conference},
\num{2575}--\num{2582}.

\bibitem[Observatorio(2024)]{observatorio2024trabajo}
Observatorio del Trabajo Digno. (2024).
\emph{Condiciones de trabajo de repartidores de plataformas
digitales en la Zona Metropolitana del Valle de México.}
Reporte técnico, México.

\bibitem[OWM(2023)]{owm2023traffic}
OpenWeather Ltd. (2023). \emph{Weather data integration for
urban traffic prediction: best practices.} Whitepaper.

\bibitem[SEMOVI(2024)]{semovi2024parque}
Secretaría de Movilidad de la Ciudad de México (SEMOVI). (2024).
\emph{Estadísticas del parque vehicular y movilidad en la
\ZMVM.} Gobierno de la \CDMX.

\bibitem[TomTom(2024a)]{tomtom2024index}
TomTom International B.V. (2024).
\emph{TomTom Traffic Index 2024: global urban congestion
report.} Amsterdam.

\bibitem[TomTom(2024b)]{tomtom2024routing}
TomTom International B.V. (2024).
\emph{TomTom Routing API: technical documentation and ETA
methodology.} Amsterdam.

\end{thebibliography}

\end{document}

% VERSIÓN MARTES 14 DE ABRIL DEL 2026